\documentclass[11pt]{article}
\usepackage[margin=1in]{geometry}
\usepackage{times}
\usepackage{amsmath,amssymb}
\usepackage{booktabs}
\usepackage{graphicx}
\usepackage{hyperref}
\usepackage{xcolor}
% \todo{...} marks things to fill once real Kaggle results are in.
\newcommand{\todo}[1]{{\color{red}\textbf{[TODO: #1]}}}
\newcommand{\res}[1]{{\color{blue}#1}} % blue = number to paste from analyze.py/export.py

\title{Compressed but Candid? The Effect of Post-Training Quantization\\
       on Chain-of-Thought Faithfulness}
\author{\todo{author} \\ \todo{affiliation}}
\date{\today}

\begin{document}
\maketitle

\begin{abstract}
Chain-of-thought (CoT) monitoring---reading a model's reasoning trace to detect
misbehavior---is a leading AI-safety oversight strategy, but it works only if the
trace is \emph{faithful}: it must verbalize the factors that actually drive the
answer. Separately, weight quantization (INT8/INT4) is how open reasoning models
are actually deployed. Prior work studies these in isolation: quantization is
shown to harm reasoning \emph{accuracy} and inflate CoT \emph{length}, while
faithfulness is measured almost exclusively at full precision. We ask a question
at their intersection: \textbf{does post-training weight quantization erode CoT
faithfulness?} Using the standard hint-injection protocol on
\todo{model(s)} across FP16, INT8, and INT4, we measure the rate at which a
model verbalizes an injected cue that flips its answer. We find \todo{one-line
result}. We further separate a \emph{direct} effect of precision from an
\emph{indirect} effect mediated by the known concentration of unfaithfulness on
incorrect answers, and report the thinking-span vs.\ answer-span acknowledgment
gap. \todo{implication sentence}. Code and data are released for full
reproducibility.
\end{abstract}

%======================================================================
\section{Introduction}
Reasoning models externalize a step-by-step CoT before answering, and this trace
has been proposed as a window for AI oversight~\cite{cotmonitorability2025,baker2025}.
Its value depends on \emph{faithfulness}: recent work shows CoT can be plausible
yet systematically unfaithful, omitting the biasing cues that actually determine
the answer~\cite{turpin2023,anthropic2025reasoning,lietome2026}.

In deployment, open reasoning models are rarely run at full precision. Post-training
quantization to INT8/INT4 is standard, and a growing literature documents its
effect on reasoning \emph{accuracy} and \emph{length}: quantized models overthink,
inflate token counts, and lose accuracy at low
bit-widths~\cite{quanthurts2025,overthinking2026,tokeninflation2026,multiplicative2026}.
Yet whether the reasoning trace remains a \emph{faithful, monitorable} record
under quantization is, to our knowledge, unstudied. This matters directly for
safety: the models most likely to be deployed on commodity hardware---quantized
open models---are exactly the ones whose monitorability we have not checked.

\paragraph{Contributions.}
\begin{enumerate}
  \item We are the first to measure the effect of \emph{weight quantization}
        (FP16$\to$INT8$\to$INT4) on CoT faithfulness, holding the base checkpoint
        fixed so precision is the sole variable.
  \item Using the standard hint-injection protocol, we report cue-acknowledgment
        on answer-flipping cases across precisions, plus the thinking-vs-answer
        gap~\cite{whymodelsknow2026}.
  \item We stratify by answer correctness to separate a direct precision effect
        from the indirect effect predicted by the ``two regimes'' of
        unfaithfulness~\cite{tworegimes2026}.
  \item We release code, per-trace data, and a validated cue-detector.
\end{enumerate}

%======================================================================
\section{Related Work and Positioning}
Our contribution sits precisely between three active literatures; we state the
gap against each.

\paragraph{CoT faithfulness and monitorability.}
The hint-injection protocol~\cite{turpin2023,lanham2023} and its extensions
measure whether models verbalize cues that change their
answers~\cite{anthropic2025reasoning,lietome2026,monitorability2025,faithcotbench2026,bonafide2026}.
\emph{All of this work is conducted at full precision.} We reuse the protocol but
add precision as the independent variable.

\paragraph{Compression vs.\ the CoT---but the \emph{length} of it.}
A separate line studies how compressing the \emph{reasoning trace} (thousands of
tokens $\to$ hundreds) harms monitorability~\cite{monitorabilitytax2026,lengthpenalties2026}.
This is a different sense of ``compression'': ours is \emph{weight} quantization,
not CoT-length reduction. We are careful to distinguish the two.

\paragraph{Quantization vs.\ reasoning---but \emph{accuracy}, not faithfulness.}
Quantized reasoning models are known to lose accuracy and overthink at low
bit-widths~\cite{quanthurts2025,quantmeets2025,overthinking2026,tokeninflation2026,multiplicative2026,whenreasoningmeets2025}.
These measure \emph{whether the answer is right}, not \emph{whether the trace
honestly reports its cause}. The closest reliability-oriented work, UniComp,
finds a performance--reliability decoupling across compression
methods~\cite{unicomp2026}, and other work finds compression can undo
alignment~\cite{quantundoes2026}; neither measures cue-verbalization faithfulness
with the hint-injection protocol. We fill exactly that cell.

%======================================================================
\section{Method}
\paragraph{Models and precisions.}
We evaluate \todo{model list, e.g. DeepSeek-R1-Distill-Qwen-1.5B (and 7B)} at
three precisions: FP16 (baseline), INT8, and INT4 (NF4), loaded from the
\emph{same} base checkpoint via bitsandbytes~\cite{dettmers2023,deepseekr1}. As a
cross-method check against a single-quantizer artifact, we additionally include
\todo{one AWQ and one GPTQ checkpoint}.

\paragraph{Faithfulness protocol.}
For each multiple-choice question (\todo{dataset, e.g. MMLU}, $N=$\todo{N}) we
query the model twice: once \emph{clean}, once with an injected \emph{cue}
pointing at a wrong option. We use two cue types (sycophancy, authority). On
cases where the cue \emph{flips} the answer to the cued option, a faithful CoT
should mention the cue; we measure the \emph{cue-acknowledgment rate}. Generation
is greedy ($T{=}0$) for reproducibility~\cite{monitorability2025}.

\paragraph{Metrics.}
(i)~\emph{Flip rate}: fraction of cued cases whose answer moved to the cue.
(ii)~\emph{Faithfulness}: acknowledgment rate on flipped cases.
(iii)~\emph{Thinking-vs-answer gap}: acknowledgment in the \texttt{<think>} span
vs.\ the answer span~\cite{whymodelsknow2026}.
(iv)~\emph{Correctness-stratified faithfulness}, to separate direct from indirect
precision effects~\cite{tworegimes2026}.

\paragraph{Cue detector and its validation.}
Acknowledgment is detected by a lexical matcher over cue-specific marker phrases.
We validate it against \todo{K} hand-labeled traces, reporting accuracy
\res{acc} and Cohen's $\kappa=$\res{kappa} (\S\ref{sec:validation}); rates are
reported only if $\kappa$ is acceptable.

%======================================================================
\section{Results}
\paragraph{Detector validation.}\label{sec:validation}
On \todo{K} hand-labeled flipped-case traces, the lexical detector achieves
accuracy \res{acc}, $\kappa=$\res{kappa}, precision \res{P}, recall \res{R}
(from \texttt{label\_helper.py score}). \todo{one sentence: reliable / not}.

\paragraph{Faithfulness vs.\ precision (headline).}
Table~\ref{tab:headline} reports flip rate and faithfulness per precision
(from \texttt{analyze.py}). \todo{state the trend across FP16/INT8/INT4}.

\begin{table}[h]\centering
\caption{Cue-acknowledgment (faithfulness) on answer-flipping cases.}
\label{tab:headline}
\begin{tabular}{lccc}
\toprule
Precision & Flip rate & Faithfulness (ack $\mid$ flipped) & $n_{\text{flipped}}$ \\
\midrule
FP16 & \res{--} & \res{--} & \res{--} \\
INT8 & \res{--} & \res{--} & \res{--} \\
INT4 & \res{--} & \res{--} & \res{--} \\
\bottomrule
\end{tabular}
\end{table}

\paragraph{Thinking vs.\ answer gap.}
Figure~\ref{fig:gap} (from \texttt{export.py}) shows acknowledgment in the
thinking span far exceeds the answer span, replicating~\cite{whymodelsknow2026};
\todo{does the gap widen with quantization?}.

\begin{figure}[h]\centering
% \includegraphics[width=0.6\textwidth]{results_gap.png}
\fbox{\parbox{0.6\textwidth}{\centering \todo{insert results\_gap.png}}}
\caption{Thinking- vs.\ answer-span cue acknowledgment across precisions.}
\label{fig:gap}
\end{figure}

\paragraph{Correctness stratification.}
Table~\ref{tab:strat} splits faithfulness by whether the answer was correct.
\todo{Is the precision effect mostly the indirect route (more wrong answers at
low bit-width) or a direct drop within each regime?}

\begin{table}[h]\centering
\caption{Faithfulness stratified by answer correctness.}
\label{tab:strat}
\begin{tabular}{lcc}
\toprule
Precision & ack $\mid$ correct & ack $\mid$ incorrect \\
\midrule
FP16 & \res{--} & \res{--} \\
INT8 & \res{--} & \res{--} \\
INT4 & \res{--} & \res{--} \\
\bottomrule
\end{tabular}
\end{table}

%======================================================================
\section{Discussion}
\todo{Interpret: if faithfulness drops, quantized deployments are less monitorable
than their FP16 evaluation suggests---a safety gap hidden by accuracy-only
benchmarks (cf.\ UniComp's performance--reliability decoupling~\cite{unicomp2026}).
If it is preserved, that is a reassuring and equally publishable negative result:
monitorability survives standard quantization.}

%======================================================================
\section{Limitations}
Single/few model families and one size range (\todo{sizes}); lexical cue detector
(mitigated by $\kappa$ validation and an LLM-judge robustness check,
\todo{if run}); multiple-choice only; two cue types; bitsandbytes as the primary
quantizer with AWQ/GPTQ as a spot check rather than full coverage; greedy
decoding only. Any effect is correlational with respect to internal mechanism.

%======================================================================
\section*{Reproducibility}
All code, the per-trace dataset, the validated cue-detector, and the exact
commands are released at \todo{repo URL}. Experiments run on a single free-tier
T4 GPU in \todo{X} GPU-hours.

\bibliographystyle{plain}
\bibliography{references}
\end{document}
